% =========================================================================
%  Anexos
% =========================================================================

\chapter{Especificação OpenAPI}
\label{anexo:openapi}

% Excerto da especificação OpenAPI ou referência ao ficheiro
% completo. Para excertos mais longos, considerar incluir como
% código com a package "listings".

A especificação completa da API está disponível no ficheiro
\texttt{openapi.yaml} no repositório do projeto
(\url{\repositorio}). A título ilustrativo apresenta-se a seguir
o excerto correspondente ao \emph{endpoint} de gestão de alertas
especificado na Sprint 1.

\begin{lstlisting}[language=, caption={Excerto de \texttt{openapi.yaml} — ignorar alerta com justificacao.}]
paths:
  /api/v1/alerts/{alertId}/ignore:
    post:
      summary: Ignorar alerta (justificacao obrigatoria)
      parameters:
        - in: path
          name: alertId
          required: true
          schema:
            type: string
      requestBody:
        required: true
        content:
          application/json:
            schema:
              type: object
              required: [reason]
              properties:
                reason: { type: string, minLength: 3 }
      responses:
        "200":
          description: Alerta ignorado
          content:
            application/json:
              schema:
                type: object
                properties:
                  id:
                    type: string
                  status:
                    type: string
                    enum: [ignorado]
                  ignoredReason:
                    type: string
        "400":
          description: Justificacao em falta ou invalida
\end{lstlisting}


\chapter{Dados de demonstração}
\label{anexo:dados}

% Descrição breve do dataset de demonstração entregue junto com
% o código (CSV ou Excel exemplo).

O dataset de demonstração inclui:

\begin{itemize}
  \item {}[N] utilizadores (um por perfil, no mínimo).
  \item {}[N] ervas aromáticas, importadas a partir do ficheiro \texttt{ervas-exemplo.csv}.
  \item {}[N] planos de cultivo (\emph{regular}, \emph{emergência} e \emph{pontual}).
  \item {}[N] lotes em diferentes estados (\emph{ativo}, \emph{concluído}, \emph{comprometido}).
  \item Histórico de medições ambientais e alertas associados.
\end{itemize}


\chapter{Distribuição do trabalho}
\label{anexo:distribuicao}

% Tabela honesta com a contribuição de cada elemento da equipa.
% Útil para a avaliação individual e para a equipa refletir.

\begin{table}[H]
\centering
\caption{Contribuição de cada elemento da equipa.}
\label{tab:distribuicao}
\renewcommand{\arraystretch}{1.4}
\begin{tabularx}{\linewidth}{l X c}
\toprule
\textbf{Elemento} & \textbf{Contribuições principais} & \textbf{\%} \\
\midrule
Aluno 1 & [Áreas de responsabilidade e principais entregas]   &      \\
Aluno 2 & [Áreas de responsabilidade e principais entregas]   &      \\
Aluno 3 & [Áreas de responsabilidade e principais entregas]   &      \\
\midrule
\textbf{Total} & & \textbf{100\%} \\
\bottomrule
\end{tabularx}
\end{table}
